\section{JASPによるベイズ分析入門}
\label{b27_jasp_intro}

\subsection{授業内容}

\subsubsection{概要}
これまで理論的に学んできたベイズ統計学を実践するために，オープンソースの統計ソフトウェアJASP（Jeffreys's Amazing Statistics Program）を導入する。JASPは直感的なGUIを備え，従来の頻度主義的手法とベイズ統計学的手法を並列的に提供する特徴を持つ。本コマでは，JASPのインストールと基本的な操作方法を学び，サンプルデータを用いた記述統計，可視化，相関分析，回帰分析を実践する。特に，従来の頻度主義的分析とベイズ的分析の結果を比較しながら，両者の出力の違いと解釈方法について理解を深める。これにより，統計ソフトウェアを用いたベイズ分析の基礎を身につけ，次回以降のより高度な分析への準備とする。

\subsubsection{コマ主題細目}
\begin{description}
	\item[JASPの導入とインターフェイス] JASPはオープンソースの統計ソフトウェアで，心理学研究のための統計分析機能と直感的なユーザーインターフェイスを提供する。まず，JASPのウェブサイト（https://jasp-stats.org/）からソフトウェアをダウンロードし，各自のPCにインストールする方法を学ぶ。続いて，JASPの基本的なインターフェイスと機能（メニュー構成，データビュー，分析ビュー，結果ビュー）について理解する。JASPの特徴として，オープンサイエンスをサポートする機能（分析結果の再現性，OSFとの連携）や，従来の頻度主義的分析とベイズ的分析を同じメニューから実行できる利便性について学ぶ。また，サンプルデータの読み込み方，外部データ（CSV，SPSSなど）のインポート方法，変数の種類の設定方法についても実習を通じて習得する。
	      \begin{flushright}
		      $\to$ \textcite{YUnaShimizu2022}のP.1--15.
	      \end{flushright}

	\item[記述統計と可視化] JASPを用いた基本的なデータ分析として，記述統計と可視化の方法を学ぶ。「Descriptives」メニューを使って，各変数の基本統計量（平均，中央値，標準偏差，四分位数など）を計算し，結果の読み方を理解する。また，ヒストグラム，箱ひげ図などのグラフを作成し，データの分布特性を視覚的に把握する方法を学ぶ。これらの基本操作を通じて，JASPのワークフローに慣れ，次のステップであるより高度な分析への準備を整える。
		\begin{flushright}
		      $\to$ \textcite{YUnaShimizu2022}のP.1--15.
	      \end{flushright}

	\item[相関分析：頻度主義的手法とベイズ的手法] JASPを用いた相関分析を，伝統的手法とベイズ的手法の両方で実施し，結果を比較する。頻度主義的アプローチでは，相関係数，$p$値，信頼区間が出力される一方，ベイズ的アプローチでは，ベイズファクター（$BF_{10}$），事後分布，確信区間が出力される違いを理解する。特に，帰無仮説（相関がない）と対立仮説（相関がある）の間のベイズファクターが，証拠の強さをどのように表現するかを学ぶ。

	\item[回帰分析のベイズ推定] JASPを用いた回帰分析を，頻度主義的アプローチとベイズ的アプローチの両方で実施する。「回帰」メニューの「線形回帰」を伝統的手法とベイズ的手法，それぞれ使用して分析を行い，結果の違いを比較する。ベイズ回帰分析では，回帰係数の事後分布，確信区間，モデル比較のためのベイズファクターが出力される。特に，従来の$p$値に基づく変数選択と，ベイズファクターに基づくモデル比較の違いを理解する。さらに，次回の内容につながる準備として，t検定や分散分析のベイズ版についても簡単に触れる。

\end{description}
\subsubsection{キーワード}
\begin{itemize}
	\item JASP（Jeffreys's Amazing Statistics Program）
	\item オープンソース統計ソフトウェア
	\item ベイズファクター（$BF_{10}$）
	\item 事後分布の可視化
\end{itemize}
\subsection{授業情報}
\paragraph{コマの展開方法}
実習
\paragraph{標準シラバスにおける位置づけ}
科目番号5；心理学統計法；(1)心理学で用いられる統計手法；J より高度な記述や推定を目指して
\subsubsection{予習・復習課題}
\paragraph{予習}
授業前にJASPをダウンロードし，各自のPCにインストールしておくこと（https://jasp-stats.org/からダウンロード可能）。インストール後，ソフトウェアが正常に起動するか確認しておくとよい。また，これまでの授業で学んだベイズ統計学の基本概念（事前分布，事後分布，ベイズファクターなど）を復習しておくこと。さらに，第\ref{b06_Corr}，\ref{b07_regression}講で学んだ相関分析と回帰分析の基本についても再確認しておくとよい。

\paragraph{復習}
授業で学んだ内容を定着させるために，以下の点について実践を通じて理解を深めておくこと：

\begin{enumerate}
\item JASPの基本的な操作（データの読み込み，変数タイプの設定，分析の実行，結果の保存）を繰り返し練習する
\item 自分で入手したデータや，JASPに内蔵されているサンプルデータを用いて，記述統計と可視化を実行してみる
\item 相関分析を頻度主義的アプローチとベイズ的アプローチの両方で実行し，結果の違いを比較検討する
\item 回帰分析についても同様に両アプローチで実行し，係数の推定値，信頼/確信区間，$p$値/ベイズファクターの違いをまとめる
\end{enumerate}

特に，JASPの操作に慣れることが重要であるため，様々なデータセットを用いて分析を繰り返し実行してみるとよい。また，JASPの公式ウェブサイトには多数のチュートリアルやガイドが掲載されているので，それらを参考にしながら学習を進めることも有効である。次回の授業では検定や分散分析のベイズ版を扱うため，基本操作を確実に身につけておくことが重要である。
